% ============================================================
% Evaluation. Organize around research questions (RQs); every
% table/figure answers exactly one RQ, and every RQ is answered.
% Numbers never stand alone — each gets an interpretation.
% ============================================================
\section{Evaluation}
\label{sec:evaluation}

% State the research questions up front; they structure the
% whole section:
% \begin{itemize}
%     \item \textbf{RQ1:} How does the proposed method compare with state-of-the-art baselines?
%     \item \textbf{RQ2:} How does each component contribute to overall performance?
%     \item \textbf{RQ3:} How sensitive is the method to parameter X?
% \end{itemize}

\subsection{Experimental Setup}
\label{subsec:setup}

% Enough detail for independent replication:
%   - Datasets: source, size, splits, preprocessing, license
%   - Baselines: which, versions/configs, why they are the right comparisons
%   - Metrics: definitions with citations; how ties/edge cases count
%   - Environment: hardware, software versions (briefly)
%   - Statistical treatment: runs, error bars, significance tests

\subsection{Overall Performance (RQ1)}
\label{subsec:results}

% Main comparison table. Use booktabs rules (no vertical lines),
% bold the best result, and order rows meaningfully (by score,
% or baselines-then-ours):
\begin{table}[htbp]
    \caption{Comparison with baselines. Best results in \textbf{bold}.}
    \label{tab:main-results}
    \centering
    \begin{tabular}{lccc}
        \toprule
        Method & Accuracy & F1 & Time (s) \\
        \midrule
        Baseline A \cite{example2023} & 85.2 & 84.1 & 1.2 \\
        Baseline B \cite{conference2022} & 87.5 & 86.8 & 2.1 \\
        \textbf{Ours} & \textbf{91.3} & \textbf{90.7} & 1.8 \\
        \bottomrule
    \end{tabular}
\end{table}

% Then interpret: not just "ours is better" but WHY, tied back
% to the mechanism in Section~\ref{sec:methodology}, and WHERE
% it wins/loses (per-dataset or per-category breakdown).

\subsection{Ablation Study (RQ2)}
\label{subsec:ablation}

% Remove/disable each component one at a time and measure the
% drop. This is where the design choices from the methodology
% get validated — or honestly refuted.

\subsection{Parameter Sensitivity (RQ3)}
\label{subsec:sensitivity}

% Vary the key hyperparameter(s) and show stability. A line
% chart works well:
% \begin{figure}[htbp]
%     \centering
%     \includegraphics[width=\linewidth]{figures/sensitivity.pdf}
%     \caption{Accuracy as a function of $\lambda$.}
%     \label{fig:sensitivity}
% \end{figure}
